\section{Conclusion}

We set out to measure whether the internal geometry of a frozen video world
model crosses an alarm threshold sooner than the model's own prediction error,
and built the experiment so the answer could come out either way. That
measurement is \pending: it needs about one GPU-hour of credit, not any missing
code. We report the apparatus, and what building it revealed.

The methodological point stands independently of the eventual verdict. A
comparison between an internal signal and an output error is unusually easy to
settle in advance without noticing: by letting the baseline's context attend to
the frames it predicts; by calibrating thresholds on the data they are scored
against; by summarising a censored latency with a statistic that is constant
under censoring; or by beating a baseline that never fires. We made the first
structurally impossible and tested it adversarially, the second impossible by
construction, the third visible through a registered secondary, and the fourth
visible through a lift check that marks an otherwise-passing verdict vacuous.

Two further findings came from building rather than from measuring. Our
occlusion condition, specified as a square covering $30\%$ of the frame, in
practice hides the disc entirely in \lgwOccFullyHidden{} of post-injection
frames --- so it probes object permanence, a regime in which the encoder and
predictor of V-JEPA 2 are known to dissociate \citep{xie2026tracking}, and
describing it as partial occlusion would have misstated what the numbers mean.
And implementing a second architecture showed that the leakage trap is
architecture-specific: VideoMAE builds its targets by patchifying whatever
tensor it is handed, reaching the same bias by an entirely different route.
Whether a loss is causal has to be re-asked, and re-tested, per architecture; it
is not a property a harness can confer.

None of these guards is sophisticated. What they share is that omitting any one
of them yields a cleaner-looking result, which is why we think they belong in
the protocol rather than in the discussion. To our knowledge the lift check in
particular is not standard in representation-based anomaly detection, though the
collapse it guards against is documented \citep{pavuluri2026asleep}.

We release the harness as a benchmark that another architecture joins in about
twenty-five lines and is then held to the same contract automatically, so that
the next claim of this shape can be checked rather than argued --- including
ours, once the run happens.

\section*{Reproducibility statement}

The protocol was frozen before any code was written and is reproduced in
Appendix~\ref{app:prereg}; every subsequent change is logged with its date and
the evidence that motivated it. The dataset is generated deterministically from
a single seed, so none is distributed. Frozen constants are pinned by a test
that fails if edited. The causality and leakage guarantees are tested, not
asserted. Every table, figure and numeric sentence in this paper is generated
from a run's result JSON by a script in the repository; with no run analysed,
the manuscript compiles with visible \textsc{pending} markers rather than
plausible values, which is why no registered number appears above.
